\section{Introduction}\label{sec:intro}

A metro line is expensive to build and almost impossible to relocate. Planners therefore need to compare many candidate extensions before committing, and they need to know who benefits from each. Network design for transit is a long-studied optimisation problem \citep{guihaire2008,farahani2013}. Exact methods and metaheuristics scale poorly once the candidate set is large, so most applications restrict the search to a few hand-drawn corridors.

Reinforcement learning (RL) offers another route. The agent builds a line station by station and learns from simulated construction which sequences give good networks. MetroGNN \citep{su2024} represents candidate regions as nodes of a graph and trains an attentive actor-critic policy on demand satisfaction. \citet{zhang2024} replace the single objective with demand, equity and accessibility, combined by Tchebycheff scalarisation. \citet{michailidis2026} show that a tabular method with fairness-based rewards can match deep RL on the Xi'an benchmark at a much smaller computational cost. \citet{holliday2025} use graph neural networks as learned heuristics inside an evolutionary algorithm for transit network design.

These studies leave three questions open for rail planning. First, the demand figure is hard to interpret. In the public Xi'an environment the percentage reward is computed per population group, and the groups overlap, so the sum over groups is not the share of total demand served (Section~\ref{sec:benchmark}). Second, the fairness measure used as a reward, one minus a Gini index, scores a line that serves nobody as perfectly fair. Third, evolutionary and learned methods are rarely compared at the same simulation budget, and rarely on all objectives. A planner choosing between a learned policy and a multi-objective genetic algorithm cannot tell from the literature which one finds better trade-offs for the same computing effort.

This paper makes the comparison explicit. We keep the Xi'an and Amsterdam benchmark environments of \citet{michailidis2026} and change how a line is scored. The demand objective is the true share of origin--destination (OD) flow served. The fairness objective is a generalised Gini welfare \citep{weymark1981} over per-group satisfaction, which is zero when nothing is served, and we also report the worst-served group directly. Accessibility is the trip-end mass newly brought within reach of the line. We then compare a weight-conditioned graph RL policy with NSGA-II \citep{deb2002}, scalarised genetic algorithms, tabular Q-learning, a greedy heuristic and random search. Every method receives about $10^4$ simulated lines, and quality is measured on the whole Pareto set by hypervolume and expected utility \citep{zitzler2003,hayes2022}.

The contributions are:
\begin{enumerate}
\item A corrected scoring protocol for the benchmark: true demand share, zero-safe fairness measures, a trip-end-based accessibility term, and objective scaling anchored on single-objective optima (Section~\ref{sec:problem}).
\item A graph RL policy trained with batched REINFORCE and a leave-one-out baseline. The policy chooses its own start cell and sees one-step lookahead features. We describe exactly what the code does, including that it uses no replay buffer (Section~\ref{sec:method}).
\item A matched-budget comparison against NSGA-II and four other baselines on two real grids, with ten seeds, exact tests and Holm correction (Sections~\ref{sec:design} and~\ref{sec:results}).
\item A leave-one-line-out test on Xi'an: remove Line~2, ask each method for a 27-station extension of Line~1, and measure how closely it recovers the real Line~2 relative to random search.
\item A paired study of cross-city warm starts between Xi'an and Amsterdam with controls that separate transferred knowledge from having fewer trainable parameters.
\end{enumerate}
We find that, at equal budget, NSGA-II gives better deployable trade-off sets than the learned policy on both grids. On Xi'an the search archives of the two are comparable, and on the larger Amsterdam grid NSGA-II is better on both measures. The policy returns almost the same line for different weight vectors. In the leave-one-line-out test, NSGA-II recovers part of the real Line~2 and the learned policy does not differ from random search. Ablations show that message passing and a scalarised prior both matter for the policy, and a warm start helps in one direction between the two cities and not in the other.
The remainder of the paper is organised as follows. Section~\ref{sec:related} reviews related work. Section~\ref{sec:problem} defines the problem and the objectives. Section~\ref{sec:method} describes the methods. Section~\ref{sec:design} gives the experimental design and Section~\ref{sec:results} the results. Section~\ref{sec:discussion} discusses limitations and Section~\ref{sec:conclusion} concludes.
